\section{Complete certificates on the ten finite residue disks}
\label{sec-fixed-finite-disks}

Use the curves, exact constants \(\alpha_i=\alpha_i^0\), and
\(f=F_0+\frac43\log_5 3\) of Section~\ref{zs-section}.
For a disk coordinate \(s\), write
\(\mathcal A=\mathbb Z_5\langle s\rangle\), the ring of restricted
series, with Gauss valuation
\(v_{\mathcal A}(\sum c_js^j)=\inf_jv_5(c_j)\).
Its coefficients tend to zero, and its derivative belongs to the same
ring. We prove congruences in this complete ring; evaluations at a few
points do not replace the entire-series assertions.
The involutions and four disk orbits are established in
Section~\ref{sec-fixed-rational-descent}.

\subsection{The representative disk above one}

Take
\begin{equation}
 z=1+5s,\qquad W(1)=8,\qquad s\in\mathbb Z_5.
 \label{ud-disk}
\end{equation}
The square-root branch is the unique one with this residue sign.
For either short elliptic model put
\begin{align}
 \psi&=\psi_9,&\phi&=x\psi_9^2-\psi_{10}\psi_8,\\
 \omega_9&=\frac{\psi_{11}\psi_8^2-\psi_7\psi_{10}^2}{4y},&
 \delta&=-\phi/\omega_9,&T&=-\phi\psi/\omega_9.
 \label{ud-division-functions}
\end{align}
These are the normalized division polynomials:
\([9](x,y)=(\phi/\psi^2,\omega_9/\psi^3)\).
Thus \(\delta=t([9]Q)/\psi_9(Q)\) and \(T=t([9]Q)\), with no
remaining division by \(\psi_9\) at nine-torsion.

\begin{lemma}[Integral functions and complete tails on the first disk]
\label{ud-integral}
On \eqref{ud-disk}, \(\phi_i,\omega_{9,i},\delta_i,\Xi\) are units
of \(\mathcal A\), \(\psi_{9,i},T_i\in5\mathcal A\), and
\(g=f/5\in\mathcal A\). If \(u=z-1\) and
\(\beta=(d\log\Xi/du)(0)\), then
\begin{equation}
 \log_5\Xi(s)-\log_5\Xi(0)=5\beta s\pmod{25\mathcal A}.
 \label{ud-log-tail}
\end{equation}
All series estimates remain valid after termwise differentiation in
\(s\).
\end{lemma}
\begin{proof}
The equation and \(2W(0)=16\) give
\(W(u)\in8+u\mathbb Z_5[[u]]\), with \(W'(0)=6\).
Both quotient coordinates are integral series and both ordinates
are units. The division recurrences divide only by \(2y\) or
\(4y\), which are units. The complete calculation below gives
\(\psi_9=0,\phi=\omega_9=1\) modulo five at the center, for both
quotients. An integral formal series with unit constant has an
integral inverse. Substitution \(u=5s\) gives a restricted series
whose degree-\(j\) coefficient is divisible by \(5^j\). It follows
that \(\phi,\omega_9,\delta\) are units on the entire disk and that
\(\psi_9,T\in5\mathcal A\). Since \(z\) is also a unit,
\(\Xi=z^{81}\delta_1/\delta_2\) is a unit.

The ratio \(\Xi(s)/\Xi(0)\) lies in \(1+5s\mathcal A\) and is
\(1+5\beta s\) modulo \(25\mathcal A\). In \(\log(1+h)\),
\(h\in5\mathcal A\), the terms \(h^j/j\), \(j\ge2\), have
valuation at least \(j-v_5(j)\ge2\), tending to infinity.
This proves \eqref{ud-log-tail} for the full tail. The constant
\(\log_5\Xi(0)\) lies in \(5\mathbb Z_5\).

Proposition~\ref{zs-tail} applies in the same complete algebra:
\[
 v_{\mathcal A}(r_jT_i^j)\ge j-2\lfloor\log_5j\rfloor,\qquad
 v_{\mathcal A}(l_jT_i^j)\ge j-\lfloor\log_5j\rfloor.
\]
Both bounds tend to infinity. The composed tails beyond degree \(N\)
have valuation at least \(\lceil(N+1)/2\rceil\). Since \(R_0\)
starts in degree four, \(R_{0i}(T_i)\in25\mathcal A\);
also \(L_i(T_i)/9\in5\mathcal A\).
As \(v_5(\alpha_i)=-1\), their squared-log terms lie in
\(5\mathcal A\). The same is true of \(\log_5\Xi\) and \(\Omega\),
so \eqref{zs-height-formula} proves \(g\in\mathcal A\).
Differentiation of a restricted integral series does not decrease
its Gauss valuation. This supplies derivative control for all the
tails, rather than only for a truncated polynomial.
\end{proof}

\begin{lemma}[The exact first reduction]
\label{ud-reduction}
On \eqref{ud-disk}, the full function satisfies
\begin{equation}
 g(s)=2s(s+1)\pmod{5\mathcal A}.
 \label{ud-polynomial}
\end{equation}
\end{lemma}
\begin{proof}
The only new local constant is computed in
\(\mathbb F_5[e]/(e^2)\), with \(z=1+e\), \(W=8+6e\).
The two points are
\((x_1,y_1)=(3+3e,1+2e)\),
\((x_2,y_2)=(1+2e,1+4e)\).
To specify the finite calculation completely, use
\begin{align}
 \psi_0&=0,&\psi_1&=1,&\psi_2&=2y,\\
 \psi_3&=3x^4+6ax^2+12bx-a^2,\\
 \psi_4&=4y(x^6+5ax^4+20bx^3-5a^2x^2-4abx-8b^2-a^3),\\
 \psi_{2m+1}&=\psi_{m+2}\psi_m^3-\psi_{m-1}\psi_{m+1}^3,\\
 \psi_{2m}&=\frac{\psi_m}{2y}
       (\psi_{m+2}\psi_{m-1}^2-\psi_{m-2}\psi_{m+1}^2).
 \label{ud-recurrences}
\end{align}
These agree with \cite[\S5.5]{SutherlandDivisionPolynomials2023}.
Every denominator used here has nonzero constant. The entire table
through the largest required index is as follows; entries are
\((\text{constant},\text{coefficient of }e)\).
\[
\begin{array}{c|cc}
 n&\psi_n(E_1)&\psi_n(E_2)\\\hline
0&(0,0)&(0,0)\\1&(1,0)&(1,0)\\2&(2,4)&(2,3)\\
3&(2,1)&(1,4)\\4&(2,2)&(4,1)\\5&(3,0)&(1,0)\\
6&(3,3)&(4,2)\\7&(3,1)&(3,4)\\8&(4,4)&(4,2)\\
9&(0,4)&(0,1)\\10&(1,0)&(1,0)\\11&(2,1)&(2,1)
\end{array}
\qquad
\begin{array}{c|cc}
 &E_1&E_2\\\hline
 \phi&(1,1)&(1,3)\\
 \omega_9&(1,4)&(1,2)\\
 \delta&(4,3)&(4,4)\\
 (\log\delta)'&2&1
\end{array}
\]
Hence \(\beta=81+2-1=2\pmod5\). By \eqref{ud-log-tail}, the
contribution to \((f(s)-f(0))/5\) from \(\log_5\Xi\) is
\(-2\beta s/81=s\) modulo \(5\mathcal A\); the \(R_0\) terms
make no contribution.

Set \(\ell_i=L_i(T_i)/9\). Direct pullback gives
\begin{equation}
 \frac{d\ell_1}{dz}=\frac{2z}{W},\qquad
 \frac{d\ell_2}{dz}=-\frac2W.
 \label{ud-pullbacks}
\end{equation}
At the center these are \(1/4,-1/4\). Since
\(R_1(0)=P_1\), \(R_2(0)=-P_2\), the exact residues
\eqref{zs-small-constants} imply
\begin{equation}
 \ell_1/5=4+4s,\quad \ell_2/5=3+s\pmod{5\mathcal A},
 \qquad 5\alpha_1=1,\quad5\alpha_2=3\pmod5.
 \label{ud-logarithm-digits}
\end{equation}
These are whole-series congruences: the pullbacks have integral
coefficients in \(u\), and integrating a term of degree \(j\ge1\)
then putting \(u=5s\) gives valuation at least
\(j+1-v_5(j+1)\ge2\), tending to infinity.
The change in the squared-log terms, divided by five, is therefore
\[
 -(4+4s)^2+3(3+s)^2-(-4^2+3\cdot3^2)=s+2s^2\pmod5.
\]
Finally \(f(0)=0\) exactly: \((1,8)\) is a rational point in the
finite chart, so the global-height identity and the full
away-from-five set of Section~\ref{sec:local-height-values} give
\(F_0(1,8)=\Omega\). Adding the logarithmic slope proves
\eqref{ud-polynomial}.
\end{proof}

\begin{theorem}[Two simple zeros on the first representative disk]
\label{ud-two-zeros}
The disk \eqref{ud-disk} has exactly two zeros of \(f\), both simple:
the point \((1,8)\), and a unique point with
\(s=4\pmod5\), equivalently \(z=21,W=3\pmod{25}\).
The four disks above \(z=\pm1\pmod5\) have exactly eight simple
zeros in total, four known rational points and four further local
points in one symmetry orbit.
\end{theorem}
\begin{proof}
The only roots of \(2s(s+1)\) in \(\mathbb F_5\) are \(0,4\),
with derivatives \(2,3\). For an integral restricted series the
Taylor estimate
\(g(x+h)=g(x)+hg'(x)\pmod{h^2\mathbb Z_5}\)
holds termwise and extends by completeness. Thus Newton iteration
gives a root in each of these two residue disks. For two roots
\(x,y\) in the same disk,
\[
 g(x)-g(y)=(x-y)(g'(y)+(x-y)H),\qquad H\in\mathbb Z_5,
\]
whose second factor is a unit; uniqueness follows. The other residues
are excluded by \eqref{ud-polynomial}. The first root is the exact
center. The derivatives at both roots are units for \(g\), so
\(f'=5g'\ne0\), giving simplicity. The expansion
\(W=8+30s\pmod{25}\) gives the second ordinate residue.
The two involutions from Section~\ref{sec-fixed-rational-descent}
transport these roots and their multiplicities to the other three
disks. This argument by itself makes no rationality assertion about
the second orbit.
\end{proof}

\subsection{The representative disk above two}

Take
\begin{equation}
 z=2+5s,\qquad W(2)=\sqrt{19}\equiv2\pmod5.
 \label{su-disk}
\end{equation}
Here \(u=z-2\) again gives integral quotient coordinates with unit
ordinates. The center reductions are
\(\psi_{9,i}=0\), \((\phi_1,\phi_2)=(4,1)\),
\((\omega_{9,1},\omega_{9,2})=(3,4)\) modulo five.
Consequently the whole unit-series argument of Lemma~\ref{ud-integral}
applies: \(T_i\in5\mathcal A\), \(\Xi\in\mathcal A^*\),
and \(g=f/5\in\mathcal A\), with the same derivative and tail control.
An integral series with unit constant is being inverted, so these
assertions cover the disk, not just its center.

\begin{lemma}[The first reduction above two]
\label{su-reduction}
On \eqref{su-disk},
\begin{equation}
 g(s)=4(s+1)^2\pmod{5\mathcal A}.
 \label{su-first}
\end{equation}
\end{lemma}
\begin{proof}
The exact lift is \(W(2)=12\pmod{25}\), with
\(W'=-138/W=1\pmod{25}\). The recurrences
\eqref{ud-recurrences} give the following constant residues and dual
values; the latter use \(z=2+e,W=2+e\).
\[
\begin{array}{c|cc}
 \bmod25&E_1&E_2\\\hline
 \psi_9&15&15\\\phi&24&1\\\omega_9&18&24\\
 \delta&7&1\\T&5&15
\end{array}
\qquad
\begin{array}{c|cc}
 (\text{value},\text{derivative})\bmod5&E_1&E_2\\\hline
 \phi&(4,0)&(1,3)\\\omega_9&(3,0)&(4,3)\\
 \delta&(2,0)&(1,1)\\(\log\delta)'&0&1
\end{array}
\]
Hence \(\beta=81/2+0-1=2\pmod5\), and the logarithmic local-height
slope in \(g\) is one. The center calculation also gives
\[
 \Xi=14,\quad\log_5\Xi=10,\quad\log_5 3=20\pmod{25},
 \qquad \ell_1/5=4,\quad\ell_2/5=2\pmod5.
\]
The differential values in \eqref{ud-pullbacks} are \(2,4\) modulo
five. The full integrated-tail estimate used in
\eqref{ud-logarithm-digits} therefore gives
\[
 \ell_1/5=4+2s,\qquad\ell_2/5=2+4s\pmod{5\mathcal A}.
\]
With \(5\alpha_i=(1,3)\) and \(\Omega/5=3\) modulo five,
substitution in \eqref{zs-height-formula} yields constant \(4\),
linear coefficient \(3\), and quadratic coefficient \(4\).
This is \eqref{su-first}. The center has not been assumed rational
or an exact zero.
\end{proof}

\begin{lemma}[Sharper tails at precision one hundred twenty-five]
\label{su-sharp-tail}
For either curve and any \(T\in5\mathcal A\),
\[
 R_0(T)\in5^4\mathcal A,\qquad L(T)-T\in5^4\mathcal A.
\]
Replacing \(L(T)/9\) by \(T/9\) in a squared-log term multiplied
by \(\alpha_i\) makes an error in \(5^4\mathcal A\).
For a five-adic unit \(q\), with \(h=q^4-1\),
\begin{equation}
 \log_5q=\tfrac14(h-h^2/2)\pmod{125}.
 \label{su-unit-log}
\end{equation}
\end{lemma}
\begin{proof}
The even series \(R_0\) starts at degree four. The valuation bound
for degree four is four; degree five is absent; degrees six and
seven have bounds four and five; all degrees at least eight have
bound \(\lceil j/2\rceil\ge4\). For \(L-T\), only degrees
\(j\ge5\) occur, and \(j-\lfloor\log_5j\rfloor\ge4\).
Both bounds tend to infinity, proving the claims for the entire
compositions. The difference of the two squares is in
\(5^5\mathcal A\), and \(v_5\alpha_i=-1\).
Finally, for every \(j\ge3\), the omitted term \(h^j/j\) has
valuation at least \(j-v_5(j)\ge3\), tending to infinity.
Changing the unit representative modulo \(125\) also changes its
logarithm only by a multiple of \(125\). This proves the exact
logarithm congruence.
\end{proof}

\begin{theorem}[No zero above two]
\label{su-no-zeros}
There is no zero of \(f\) in \eqref{su-disk}, or in any of the four
disks above \(z=\pm2\pmod5\).
\end{theorem}
\begin{proof}
The only possible root residue in \eqref{su-first} is \(s=4\).
At its exact center \(z=22\), the selected branch has
\(W=32\pmod{125}\). This satisfies the sextic congruence with
unit derivative. Unit denominators in \eqref{ud-division-functions}
therefore make evaluation at this residue exact modulo \(125\).
The same recurrences give
\[
 \Xi=24,\quad T_1=90,\quad T_2=10,\quad
 \log_5\Xi=100,\quad\log_5 3=95\pmod{125}.
\]
The logarithms have the full-tail justification
\eqref{su-unit-log}. Lemma~\ref{su-sharp-tail} allows omission of
the \(R_0\) terms and replacement of \(\ell_i\) by \(T_i/9\)
at this precision. Thus \(\ell_1/5=2,\ell_2/5=3\pmod{25}\).
The exact constants \eqref{zs-small-constants} give
\(5\alpha_1=16,5\alpha_2=3\pmod{25}\); equivalently these follow
from \(21/9^2=16\), \(2/22^2=3\) modulo \(25\).
All used digits lie within the previously proved nineteen-digit
absolute error for \(\alpha_i\). Substitution now gives
\begin{equation}
 g(4)=-\frac2{81}\,20-16\cdot2^2+3\cdot3^2
                     +\frac43\,19=15\pmod{25}.
 \label{su-obstruction}
\end{equation}
Every denominator is a five-adic unit. The positive \(4/3\) sign
comes from \(f=F_0-\Omega\).

Since \(g\in\mathcal A\) and \(g'(4)=0\pmod5\), its entire
Taylor expansion gives
\[
 g(4+5t)=g(4)+5tg'(4)=15\pmod{25\mathbb Z_5\langle t\rangle}.
\]
Every higher term is a multiple of \(25\). Completeness justifies
this expansion for the full restricted series and every
\(t\in\mathbb Z_5\). It excludes the only possible root residue;
the other residues were already excluded by \eqref{su-first}.
The involutions transport this exclusion to the three companion
disks.
\end{proof}

\subsection{The zero fiber and its removable pole}

Take
\begin{equation}
 z=5s,\qquad W(0)=W_0=\sqrt{-9}\equiv1\pmod5.
 \label{zd-disk}
\end{equation}
Its center is not being assumed rational or a zero of the height
function.

\begin{lemma}[An integral even removable factor]
\label{zd-removable}
On \eqref{zd-disk}, \(\Xi\) extends to an integral even unit,
\(T_i\in5\mathcal A\), and \(g=f/5\in\mathcal A\) is even.
Its central value is obtained from
\begin{equation}
 \Xi(0)=\delta_1(-9/4,W_0/8)(-W_0/2)^{81}.
 \label{zd-center-xi}
\end{equation}
\end{lemma}
\begin{proof}
For a good integral short model at five the integral formal group
law and the leading coefficient of the odd division polynomial give
\[
 \frac{t([9]Q)}t=9+t\mathbb Z_5[[t]],\qquad
 t^{80}\psi_9(Q)=9+t\mathbb Z_5[[t]].
\]
Indeed \(\psi_9\) is an integral polynomial in \(x\) of degree
forty with leading coefficient nine, and
\(x=t^{-2}(1+t\mathbb Z_5[[t]])\).
Both constants are units at five. Hence
\begin{equation}
 d(t):=\frac{\delta(Q)}{t^{81}}
   =\frac{t([9]Q)/t}{t^{80}\psi_9(Q)}
       \in1+t\mathbb Z_5[[t]].
 \label{zd-formal-ratio}
\end{equation}
This series is even, since \(\delta\) and \(t\) both change sign
under negation. This step uses the rational formal group, not an
unproved integrality assertion for zero-slope sigma.

The exact original parameter of the second quotient is
\[
 \tau(z)=t(R_2)=\frac{2z(11z^2-3)}{3W(z)}.
\]
The ordinate is an even integral series with unit constant.
Thus \(\tau\) is odd and \(\tau/z\) is a unit with constant
\(-2/W_0\). The second quotient lies in the formal group throughout
the disk, including \(O\) at its center. Consequently
\[
 \Xi(z)=\left(\frac z{\tau(z)}\right)^{81}
          \frac{\delta_1(R_1(z))}{d_2(\tau(z))}.
\]
For the first quotient the center has
\(\psi_9=0,\phi=4,\omega_9=2\pmod5\). Its integral series
therefore gives \(\delta_1\) a unit and \(T_1\in5\mathcal A\),
as in Lemma~\ref{ud-integral}. The second has
\(\tau\in5\mathcal A\), so formal multiplication gives
\(T_2\in5\mathcal A\). The displayed formula for \(\Xi\) is an
even integral unit and gives \eqref{zd-center-xi}.

All composition bounds of Lemma~\ref{ud-integral} now apply to both
quotients, giving \(g\in\mathcal A\). Moreover
\(\Xi(z)/\Xi(0)\in1+z^2\mathbb Z_5[[z]]\), so after \(z=5s\)
its logarithmic change belongs to \(25s^2\mathcal A\).
The logarithm converges in the Gauss norm with input in
\(25\mathcal A\). The \(R_0\) terms have their established full
tails. Finally, negating \(z\) fixes \(R_1\) and negates \(R_2\);
the height and squared logarithms are unchanged. Thus \(g\) is even.
\end{proof}

\begin{theorem}[No zero in either zero-fiber disk]
\label{zd-no-zeros}
On \eqref{zd-disk},
\begin{equation}
 g(s)=2s^2\pmod{5\mathcal A},\qquad g(0)=20\pmod{25}.
 \label{zd-obstruction}
\end{equation}
Consequently neither zero-fiber disk contains a zero of \(f\).
\end{theorem}
\begin{proof}
The chosen central square root has Hensel residues
\(W_0=21\pmod{25}\), \(W_0=46\pmod{125}\), with unit derivative.
Evaluate \((-9/4,W_0/8)\) by \eqref{ud-recurrences} and the
cancelled formula \eqref{zd-center-xi}. The exact finite inputs are
\[
\begin{array}{c|rr}
 &\bmod25&\bmod125\\\hline
 \Xi(0)&16&66\\
 \log_5\Xi(0)&15&15\\
 T_1(0)&20&95\\
 \log_5 3&20&95
\end{array}
\]
They give \(\ell_1(0)/5=1\pmod5\) and
\(\ell_1(0)/5=16\pmod{25}\). Every rational denominator is a unit.
The logarithms use \eqref{su-unit-log}; the entire error estimates
of Lemma~\ref{su-sharp-tail} permit omission of \(R_0\) and
replacement of \(\ell_1\) by \(T_1/9\) at precision \(125\).
At the center \(R_2=O\), so \(T_2=\ell_2=0\) exactly.
With \eqref{zs-small-constants}, first precision gives
\[
 g(0)=-\frac2{81}\,3-1^2+\frac43\,4=0\pmod5.
\]
The pullback differentials \eqref{ud-pullbacks}, integrated with
their full coefficient bounds, give
\[
 \ell_1/5=1,\qquad \ell_2/5=(-2/W_0)s=3s\pmod{5\mathcal A}.
\]
The first differential has no constant term in \(z\); all terms
after integration and substitution, apart from the stated linear
term of the second logarithm, lie in \(25\mathcal A\).
The \(\Xi\) variation and \(R_0\) tails contribute no nonconstant
term modulo five to \(g\), by Lemma~\ref{zd-removable}.
The squared-log contribution is therefore \(3(3s)^2=2s^2\).
This proves the first congruence in \eqref{zd-obstruction}.

At second precision the same table instead gives
\[
 g(0)=-\frac2{81}\,3-16\cdot16^2+\frac43\,19=20\pmod{25}.
\]
This is a direct value of the cancelled analytic formula; no
rational-point identity has been applied to the algebraic center.
Any zero would have \(s=0\pmod5\). Since \(g\) is even and
integral restricted,
\[
 g(5t)=g(0)=20\pmod{25\mathbb Z_5\langle t\rangle}.
\]
All nonconstant even terms are divisible by \(25\), and
completeness makes this an entire-series congruence. It excludes
the only possible root residue. The hyperelliptic involution gives
the same exclusion on the second zero-fiber disk.
\end{proof}

\paragraph{Finite evidence and proof boundary.}
The three exact input certificates are replayed, respectively, by
\texttt{next\_replay\_unit\_disk.py},
\texttt{next\_replay\_second\_unit\_disk.py}, and
\texttt{next\_replay\_zero\_fiber.py}. Their standard-library
calculations include the normalized division recurrences, dual-number
derivatives, Hensel square-root residues, exact unit logarithms and
global constant residues. The canonical results, source hashes and
check commands accompany the geometry inventory. The analytic proofs
above separately establish the complete tails, unit denominators,
derivative estimates, and exclusion or uniqueness on each whole disk.
These are ordinary proofs with exact finite verification; no new Lean
formalization of local heights or analytic zero counting is claimed.
The separate infinity certificate supplies the last two disks, and
the rational sieve decides the rationality of the additional first-disk
orbit. Those subsequent conclusions do not change the individual
local zero counts proved here.
